\section{Discussion and conclusion}
\label{sec:discussion}

\textbf{The horizon of the reward selects which behaviour pays.} Both rewards paid for
non-specific praise in the middle of training, but only the $\Kz$ reward, which stops at the reply,
kept paying for it: at iteration 9 it still scored praising replies $0.16$ within-group standard
deviations above their siblings, where the $\Kf$ reward, which reads five turns further, paid
$0.04$ (Appendix~\ref{app:premium}). The $\Kz$ policy ends up praising in $0.41$ of its turns and
in reply to a third of the patient's sustain talk, and late in the session its patients voice less
change talk than the Base's. What the $\Kf$ run learned instead is to keep the patient's change
talk going: its patients voice change talk again after $97\%$ of their change talk, against $80\%$
under $\Kz$. Part of that is what MI prescribes, reflecting change talk; part is not, since the
$\Kf$ policy meets sustain talk with persuasion, and the held-out judge sees more praise in it than
the training oracle does. Extending the horizon does not remove reward hacking; it changes which
behaviour pays. How the horizon produces that difference inside the update we could not isolate:
three candidate mechanisms were testable and none is confirmed (Appendix~\ref{app:mechanism});
what remains is a small, repeated difference in which sibling the group prefers, compounded over
ten iterations of a self-training loop, and the process codes show what it amounts to.
\label{sec:mechanism}

\textbf{Conclusion.} For a group-relative optimiser in multi-turn dialogue, what the judge is
shown is a first-order design choice. Scoring the continuation rather than the turn raised what
the same optimiser extracted from the same oracle, and it decided what kind of therapist the
policy became: one that praises whatever the patient said, or one that reflects the patient's
change talk and hears more of it, though it still persuades when the patient resists. We recommend
reporting the reward horizon as a hyperparameter, coding the process as well as the outcome, and
keeping a second judge with headroom (Appendix~\ref{app:saturation}).
